% ============================================================
% Section: Problem Formulation
% ============================================================
\section{Problem Formulation}
\label{sec:problem}

We consider the steady-state two-group neutron diffusion problem on a square domain.
The goal is to construct analytical functions $\varphi_1(x,y)$ and $\varphi_2(x,y)$ that simultaneously satisfy the governing PDEs and all boundary conditions.

% ------------------------------------------------------------
\subsection{Governing Equations}
\label{sec:governing}

The two-group neutron diffusion equations couple the fast-group flux $\varphi_1$ and the thermal-group flux $\varphi_2$ as follows.

\paragraph{Fast-group equation.}
\begin{equation}
\label{eq:fast_group}
-D_1 \nabla^2 \varphi_1 + \Sigma_r \, \varphi_1
= \nu \Sigma_{f1} \, \varphi_1 + \nu \Sigma_{f2} \, \varphi_2,
\end{equation}

\paragraph{Thermal-group equation.}
\begin{equation}
\label{eq:thermal_group}
-D_2 \nabla^2 \varphi_2 + \Sigma_{a2} \, \varphi_2
= \Sigma_{1\to2} \, \varphi_1,
\end{equation}

\noindent where the Laplacian operator is defined as
$\nabla^2 \varphi = \frac{\partial^2 \varphi}{\partial x^2} + \frac{\partial^2 \varphi}{\partial y^2}$.

% ------------------------------------------------------------
\subsection{Physical Parameters}
\label{sec:parameters}

Table~\ref{tab:params} summarizes the physical constants used throughout this work.

\begin{table}[h]
\centering
\caption{Physical parameters for the two-group neutron diffusion problem.}
\label{tab:params}
\begin{tabular}{@{}lll@{}}
\toprule
Symbol & Description & Value \\
\midrule
$D_1$               & Fast-group diffusion coefficient        & $1.0$   \\
$D_2$               & Thermal-group diffusion coefficient     & $0.5$   \\
$\Sigma_r$          & Fast-group removal cross-section        & $0.02$  \\
$\Sigma_{a2}$       & Thermal-group absorption cross-section  & $0.1$   \\
$\nu$               & Average neutrons per fission            & $2.5$   \\
$\Sigma_{f1}$       & Fast-group fission cross-section        & $0.005$ \\
$\Sigma_{f2}$       & Thermal-group fission cross-section     & $0.1$   \\
$\Sigma_{1\to2}$    & Group transfer cross-section            & $0.015$ \\
\bottomrule
\end{tabular}
\end{table}

% ------------------------------------------------------------
\subsection{Domain and Boundary Conditions}
\label{sec:boundary}

The spatial domain is a unit square centered at the origin:
\begin{equation}
\label{eq:domain}
\Omega = [-0.5,\, 0.5]^2 \subset \mathbb{R}^2.
\end{equation}

All four boundaries carry Neumann (current) conditions. Specifically, the left boundary carries a non-zero source, while the remaining three boundaries are zero-flux:

\paragraph{Left boundary ($x = -0.5$, non-zero Neumann).}
\begin{equation}
\label{eq:bc_left}
-D_1 \frac{\partial \varphi_1}{\partial x}\bigg|_{x=-0.5} = y,
\qquad
-D_2 \frac{\partial \varphi_2}{\partial x}\bigg|_{x=-0.5} = y.
\end{equation}

\paragraph{Right boundary ($x = 0.5$, zero Neumann).}
\begin{equation}
\label{eq:bc_right}
-D_1 \frac{\partial \varphi_1}{\partial x}\bigg|_{x=0.5} = 0,
\qquad
-D_2 \frac{\partial \varphi_2}{\partial x}\bigg|_{x=0.5} = 0.
\end{equation}

\paragraph{Top boundary ($y = 0.5$, zero Neumann).}
\begin{equation}
\label{eq:bc_top}
-D_1 \frac{\partial \varphi_1}{\partial y}\bigg|_{y=0.5} = 0,
\qquad
-D_2 \frac{\partial \varphi_2}{\partial y}\bigg|_{y=0.5} = 0.
\end{equation}

\paragraph{Bottom boundary ($y = -0.5$, zero Neumann).}
\begin{equation}
\label{eq:bc_bottom}
-D_1 \frac{\partial \varphi_1}{\partial y}\bigg|_{y=-0.5} = 0,
\qquad
-D_2 \frac{\partial \varphi_2}{\partial y}\bigg|_{y=-0.5} = 0.
\end{equation}

% ------------------------------------------------------------
\subsection{Evaluation Criteria}
\label{sec:evaluation}

A candidate solution $(\varphi_1, \varphi_2)$ is assessed by its PDE residuals and boundary condition residuals.

\paragraph{PDE residuals.}
At each interior test point $(x, y) \in \Omega$, we compute:
\begin{align}
\label{eq:res_fast}
\mathrm{Res}_1(x,y) &= -D_1 \left(\frac{\partial^2 \varphi_1}{\partial x^2} + \frac{\partial^2 \varphi_1}{\partial y^2}\right) + \Sigma_r \, \varphi_1 - \nu \Sigma_{f1} \, \varphi_1 - \nu \Sigma_{f2} \, \varphi_2, \\
\label{eq:res_thermal}
\mathrm{Res}_2(x,y) &= -D_2 \left(\frac{\partial^2 \varphi_2}{\partial x^2} + \frac{\partial^2 \varphi_2}{\partial y^2}\right) + \Sigma_{a2} \, \varphi_2 - \Sigma_{1\to2} \, \varphi_1.
\end{align}

The test points used for verification are:
\begin{equation}
\label{eq:test_points}
\mathcal{T} = \{(0,0),\; (0.2, 0.2),\; (-0.2, -0.3),\; (0.4, -0.4)\}.
\end{equation}

\paragraph{Boundary condition residuals.}
For each boundary segment $\Gamma_i$, we evaluate the deviation from the prescribed Neumann data. The combined score aggregates both PDE and BC residuals as root-mean-square (RMS) values.

\paragraph{Combined score.}
The overall quality metric is defined as:
\begin{equation}
\label{eq:combined_score}
S = \frac{1}{1 + \bar{R}_{\mathrm{PDE}} + \bar{R}_{\mathrm{BC}}},
\end{equation}
where $\bar{R}_{\mathrm{PDE}}$ and $\bar{R}_{\mathrm{BC}}$ denote the root-mean-square (RMS) residuals over the interior test points and boundary evaluation points, respectively. The score $S \in (0, 1]$, with $S = 1$ indicating a perfect solution. Programs that fail to execute or produce invalid outputs receive $S = 0$.
